\documentclass[11pt]{amsart}
\usepackage{geometry}                % See geometry.pdf to learn the layout options. There are lots.
\geometry{letterpaper}                   % ... or a4paper or a5paper or ... 
%\geometry{landscape}                % Activate for for rotated page geometry
%\usepackage[parfill]{parskip}    % Activate to begin paragraphs with an empty line rather than an indent
\usepackage{graphicx}
\usepackage{amssymb}
\usepackage{epstopdf}
\DeclareGraphicsRule{.tif}{png}{.png}{`convert #1 `dirname #1`/`basename #1 .tif`.png}

\usepackage{enumerate}
\usepackage{nicefrac}

\usepackage{mathrsfs}

\newcommand{\R}{\mathbb{R}}
\newcommand{\tstD}{\mathscr{D}}
\newcommand{\semP}{\mathscr{P}}
\newcommand{\eps}{\varepsilon}
\newcommand{\supp}{\mathrm{supp}}
\newcommand{\exd}{\mathrm{d}}
\newcommand{\loc}{\mathrm{loc}}

\title{Math 580 Assignment 3}
\author{Due Monday November 11}
\date{Fall 2013}                                           % Activate to display a given date or no date

\begin{document}
\maketitle

\begin{enumerate}[1.]
\item
Let $Q=(0,1)^n$ and let $Q_h=(h,1-h)^n$.
For $h>0$ small, define the trace map $\gamma_h:C^1(Q)\to C(\partial Q_h)$ by $\gamma_h v = v|_{\partial Q_h}$.
\begin{enumerate}[a)]
\item
Prove that $\gamma_h$ can be uniquely extended to a bounded map $\gamma_h:H^1(Q)\to L^2(\partial Q_h)$.
\item
Make sense of the boundary trace $\displaystyle\gamma_0u=\lim_{h\to0}\gamma_hu$ in $L^2(\partial Q)$ for $u\in H^1(Q)$.
\item
Show that $\gamma_0u=0$ for $u\in H^1_0(Q)$.
\item
Let $u\in H^1_0(Q)$ and let $u$  be continuous at $0$.
Show that $u(0)=0$.
\end{enumerate}
\item
Let $\Omega\subset\R^n$ be a domain, and let $W^{1,1}_\loc(\Omega)$ be the set of locally integrable functions whose (weak) derivatives are locally integrable (that is, in $L^1_\loc(\Omega)$).
\begin{enumerate}[a)]
\item% (Product rule)
Show that if $u,v\in W^{1,1}_\loc(\Omega)$ and $uv,u\partial_iv+v\partial_iu\in L^{1}_\loc(\Omega)$, 
then $uv\in W^{1,1}_\loc(\Omega)$ and $\partial_i(uv)=u\partial_iv+v\partial_iu$.
\item% (Coordinate change)
Let $\phi:\Omega\to\Omega'$ be a $C^1$-diffeomorphism between $\Omega$ and $\Omega'$.
Show that if $u\in W^{1,1}_\loc(\Omega')$ then $v=u\circ\phi\in W^{1,1}_\loc(\Omega)$ and
$\partial_iv(x) = \sum_j\partial_i\phi_j(x)(\partial_ju)(\phi(x))$, where $\phi_j$ is the $j$-th component of $\phi$, 
and $(\partial_ju)(\phi(x))$ is the evaluation of $\partial_ju$ at the point $\phi(x)$.
\item% (Chain rule)
Let $f\in C^1(\R)$ with both $f$ and $f'$ bounded, and let $u\in W^{1,1}_\loc(\Omega)$.
Prove that $f\circ u \in W^{1,1}_\loc(\Omega)$ and that $\partial_i (f\circ u) = (f'\circ u) \partial_i u$.
\item% (Chain rule with a piecewise smooth function)
Let $u\in W^{1,1}_\loc(\Omega)$ and let $u^+=\max\{u,0\}$ and $u^-=\min\{u,0\}$ pointwise.
Prove that $\partial_iu^+=\theta(u) \partial_iu$ and $\partial_iu^-=\theta(-u) \partial_iu$ {\em a.e.},
where $\theta$ is the Heaviside step function.
In particular, show that $|u|\in W^{1,p}(\Omega)$ if $u\in W^{1,p}(\Omega)$.
\end{enumerate}
\item %(Riesz representation theorem)
Let $H$ be a (real) Hilbert space, and let $H'$ be its dual, 
defined as the space of continuous linear functionals on $H$.
Let us denote the inner product of $H$ by $\langle\cdot,\cdot\rangle$.
Observe that any $y\in H$ defines an element $f\in H'$ by $f(x)= \langle y,x\rangle$ for $x\in H$.
This defines a map $J:H\to H'$.
The {\em Riesz representation theorem} (for Hilbert spaces)\footnote{There is another result called Riesz representation theorem that is about representing linear functionals on a space of continuous functions as measures.}
states that $J$ is invertible, that is, any continuous linear functional on $H$ can be realized through the inner product with an element of $H$.
We would like to prove this theorem by using a variational method.
Let $f\in H'$, and let 
$$
E(x) = \langle x,x\rangle - 2f(x),
\qquad x\in H,
$$
and consider the problem of finding a minimizer of $E$ over $H$.
\begin{enumerate}[a)]
\item Show that a minimizing sequence for $E$ exists and is Cauchy in $H$.
\item Demonstrate that the limit minimizes $E$ over $H$.
\item Denoting by $y\in H$ the minimizer, show that $\langle y,x\rangle=f(x)$ for all $x\in H$.
\item Finally, show that $y$ depends continuously on $f$.
\end{enumerate}
\item
Let $\Omega\subset\R^n$ be a bounded smooth  domain, and
consider the bilinear form
$$
a(u,v) = \int_\Omega ( a_{ij}\partial_iu\partial_jv + cuv),
$$
where the repeated indices are summer over, and the coefficients $a_{ij}$ and $c$ are smooth functions on $\overline{\Omega}$,
with $a_{ij}$ satisfying the uniform ellipticity condition
$$
a_{ij}(x)\xi_i\xi_j\geq\lambda |\xi|^2,
\qquad \xi\in\R^n,
\quad x\in\overline\Omega,
$$
for some constant $\lambda>0$.
\begin{enumerate}[a)]
\item
Show that the mapping $A:H^1_0(\Omega)\to[H^1_0(\Omega)]'$, defined by $\langle Au,v\rangle=a(u,v)$, is bounded, where $\langle\cdot,\cdot\rangle$ is the duality pairing between $[H^1_0(\Omega)]'$ and $H^1_0(\Omega)$.
\item
Show that if $c\geq0$ then 
\begin{equation*}
\langle Au,u\rangle \geq \alpha\|u\|_{H^1}^2,
\qquad u\in H^1_0(\Omega),
\end{equation*}
for some constant $\alpha>0$.
Show also that the inequality is still true (with possibly different $\alpha>0$) if $c$ is slightly negative.
\item
Supposing that $c\geq0$, show that given $f\in L^2(\Omega)$, there exists a unique function $u\in H^1_0(\Omega)$ satisfying $a(u,v)=\int_\Omega fv$ for all $v\in H^1_0(\Omega)$.
\item
Suppose that $u\in H^1_0(\Omega)$ is sufficiently smooth and satisfies $a(u,v)=\int_\Omega fv$ for all $v\in H^1_0(\Omega)$.
What differential equation does $u$ satisfy in $\Omega$?
Is $u=0$ on $\partial\Omega$?
In the language of variational methods,
this is an {\em essential} boundary condition because it is incorporated into the space $H^1_0(\Omega)$.
\end{enumerate}
\item
Let $a$ be the bilinear form as in the preceding question.
\begin{enumerate}[a)]
\item
Show that the mapping $A:H^1(\Omega)\to[H^1(\Omega)]'$, defined by $\langle Au,v\rangle=a(u,v)$, is bounded, where $\langle\cdot,\cdot\rangle$ is the duality pairing between $[H^1(\Omega)]'$ and $H^1(\Omega)$.
\item
Show that if $c>0$ in $\overline\Omega$, then
\begin{equation*}
\langle Au,u\rangle \geq \alpha\|u\|_{H^1}^2,
\qquad u\in H^1(\Omega),
\end{equation*}
for some constant $\alpha>0$.
\item
Supposing that the condition in b) holds, show that given $f\in L^2(\Omega)$, there exists a unique function $u\in H^1(\Omega)$ satisfying $a(u,v)=\int_\Omega fv$ for all $v\in H^1(\Omega)$.
\item
Suppose that $u\in H^1(\Omega)$ is sufficiently smooth and satisfies $a(u,v)=\int_\Omega fv$ for all $v\in H^1(\Omega)$.
What differential equation does $u$ satisfy in $\Omega$?
What boundary condition does $u$ satisfy?
This is a {\em natural} boundary condition because it arises from the equation $u$ has to satisfy in the weak sense.
\end{enumerate}
\item
In the context of the preceding question, assuming $c\equiv0$, prove that 
there exists a function $u\in H^1(\Omega)$ satisfying $a(u,v)=\int_\Omega fv$ for all $v\in H^1(\Omega)$ if and only if $\int f = 0$.
Moreover, show that such a function is unique up to addition of a constant.
\end{enumerate}

\end{document}  












